\documentclass[12pt,a4paper]{article}

\usepackage[utf8]{inputenc}
\usepackage[T1]{fontenc}
\usepackage{amsmath,amssymb,amsfonts}
\usepackage{geometry}
\usepackage{setspace}
\usepackage{hyperref}
\usepackage{cleveref}

\geometry{margin=2.5cm}
\onehalfspacing

\title{\textbf{Literature Survey: Expressive Bounds of State Space Models}\\[6pt]
\large Anchoring Fuchs's Cross-Architecture Observer Test}
\author{Rupert Giles\\[4pt]
\textit{Lab Research Librarian}}
\date{March 2026}

\begin{document}
\maketitle

\begin{abstract}
This survey documents the literature regarding the computational and expressive limits of State Space Models (SSMs), specifically in comparison to the well-documented $\mathsf{TC}^0$ bounds of Transformer architectures. This survey provides the theoretical anchoring required for Fuchs's "Cross-Architecture Observer Test," which aims to determine if architectural limits consistently map to specific structural fractures in output distributions.
\end{abstract}

\section{Introduction}

Fuchs has recently filed the "Cross-Architecture Observer Test" to empirically adjudicate the debate between Aaronson's Algorithmic Collapse and Wolfram's Observer-Dependent Physics. The test proposes replacing the Transformer architecture with an SSM (e.g., Mamba) to see if the semantic noise ($\Delta$) observed under narrative framing takes on a new, characteristic structure correlated with the new architecture's specific heuristic limits.

For this test to be meaningful, the baseline expressive limits of SSMs must be formally understood. Do SSMs bypass the $\mathsf{TC}^0$ bounds that cripple Transformers in $O(1)$ depth, or do they share similar algorithmic ceilings despite their recurrent formulation?

\section{Literature on SSM Expressivity}

The following papers rigorously establish the computational limits of State Space Models.

\vspace{1em}
\noindent\textbf{1. The Illusion of State in State-Space Models}\\
\textit{Merrill, W. et al. (2024). arXiv:2404.08819.}
\begin{itemize}
    \item \textbf{Relevance:} Directly addresses whether the recurrent formulation of SSMs grants them superior state-tracking capabilities compared to Transformers.
    \item \textbf{Key Finding:} The expressive power of SSMs is limited very similarly to transformers; SSMs cannot express computation outside the complexity class $\mathsf{TC}^0$. They fail at simple state-tracking problems like permutation composition. This paper formally anchors the expectation that SSMs will also fail the Rosencrantz structural grid test, as both architectures share the $\mathsf{TC}^0$ ceiling.
\end{itemize}

\vspace{1em}
\noindent\textbf{2. The Expressive Capacity of State Space Models: A Formal Language Perspective}\\
\textit{Sarrof, Y. et al. (2024). arXiv:2405.17394.}
\begin{itemize}
    \item \textbf{Relevance:} Provides a formal language perspective comparing SSMs and Transformers.
    \item \textbf{Key Finding:} While SSMs and Transformers have distinct strengths (e.g., SSMs handle star-free state tracking exactly), current SSM design choices still impose strict limits on their expressive power. This implies that the specific *type* of algorithmic failure an SSM experiences under combinatorial stress may differ from a Transformer, directly supporting Wolfram's prediction of "Observer-Dependent Physics" producing uniquely structured deviation distributions ($\Delta_{SSM} \neq \Delta_{Transformer}$).
\end{itemize}

\section{Conclusion}

The literature confirms that while SSMs differ architecturally from Transformers, they share the fundamental $\mathsf{TC}^0$ complexity bound and struggle with dynamic state tracking. This theoretical grounding validates Fuchs's experimental design: comparing two distinct but formally bounded architectures to determine if their specific structural limits produce characteristic, observer-dependent semantic fractures when faced with \#P-hard ground truths.

\end{document}
